\documentclass[11pt]{article}
\usepackage[margin=1in]{geometry}
\usepackage{amsmath,amssymb,amsthm}
\usepackage{graphicx}
\usepackage[hidelinks]{hyperref}

\newtheorem{theorem}{Theorem}
\newtheorem{corollary}{Corollary}
\newtheorem{definition}{Definition}
\newtheorem{remark}{Remark}

\newcommand{\Lip}{\operatorname{Lip}_0}
\newcommand{\PNA}{\operatorname{PNA}}
\newcommand{\SNA}{\operatorname{SNA}}

\title{A Norm-One Extension Lifting Principle for \(c_0\) in Pointwise Norm-Attaining Lipschitz Functions}
\author{agent\_lane\_04}
\date{June 29, 2026}

\begin{document}
\maketitle

\section*{Status}

This packet gives a likely valid partial result toward Question 4.1 of
Choi--Jung--Lee--Roldan \cite{CJLR231200393}.  The source asks whether
\(c_0\) is isometrically contained in \(\PNA(M)\) for every infinite pointed
metric space \(M\).  The theorem below proves an exact permanence principle:
the answer is positive whenever a positive subspace \(K\subseteq M\) admits a
norm-one linear Lipschitz extension operator into \(M\).  In particular, it is
positive when such a \(K\) is a pointed 1-Lipschitz retract of \(M\).

\section*{Source Question}

Question 4.1 in the source asks:
\begin{quote}
Is it true for every infinite pointed metric space \(M\) that \(c_0\) is
isometrically contained in \(\PNA(M)\)?
\end{quote}

\begin{center}
\includegraphics[width=0.86\textwidth]{figures/open_problem_crop.png}
\end{center}

The source proves that every infinite \(M\) contains a metric subspace
\(M_0\subseteq M\) such that \(c_0\) is isometrically contained in
\(\PNA(M_0)\).  The remaining obstruction is the subspace-to-ambient lifting
step: scalar norm-preserving Lipschitz extensions are available one function
at a time, but they are not automatically linear on an entire \(c_0\)-copy.

\section*{Definitions}

Let \(M\) be a pointed metric space.  The Banach space \(\Lip(M)\) consists
of all scalar Lipschitz functions \(f:M\to\mathbb R\) with \(f(0)=0\), equipped
with the Lipschitz norm
\[
  \|f\|_L=\sup_{p\ne q}\frac{|f(p)-f(q)|}{d(p,q)}.
\]
Following the source, \(f\in \Lip(M)\) is pointwise norm-attaining if there is
some \(p\in M\) such that
\[
  \sup_{q\in M\setminus\{p\}}
  \frac{|f(p)-f(q)|}{d(p,q)}=\|f\|_L.
\]
The set of such functions is denoted by \(\PNA(M)\).

\section*{Proof Intuition}

The source's subspace theorem almost solves the problem, except that
Lipschitz extension from \(M_0\) to \(M\) is usually nonlinear.  A nonlinear
McShane extension preserves each individual norm, but it does not preserve the
linear \(c_0\)-structure.

The obstruction disappears exactly when the relevant subspace has a
simultaneous linear extension operator of norm one.  If \(J:c_0\to\Lip(K)\)
is an isometric copy consisting of pointwise norm-attaining functions and
\(E:\Lip(K)\to\Lip(M)\) is a norm-one extension operator, then \(EJ\) is still
linear.  It is contractive because \(E\) has norm one, and it is isometric
because restricting \(EJ(a)\) back to \(K\) recovers \(J(a)\).  Finally, the
same point in \(K\) that witnesses pointwise norm-attainment for \(J(a)\) also
witnesses it for \(EJ(a)\), since the supremum over \(M\) contains the
supremum over \(K\).

\section*{Main Theorem}

\begin{theorem}[Norm-one extension lifting]
\label{thm:lifting}
Let \(K\subseteq M\) be pointed metric spaces with the same distinguished
point.  Suppose there is a linear operator
\[
  E:\Lip(K)\longrightarrow \Lip(M)
\]
such that \(\|E\|=1\) and \((Ef)|_K=f\) for every \(f\in\Lip(K)\).  If
\(\PNA(K)\) contains an isometric copy of \(c_0\), then \(\PNA(M)\) contains
an isometric copy of \(c_0\).
\end{theorem}

\begin{proof}
Let \(J:c_0\to \Lip(K)\) be a linear isometry such that
\[
  J(c_0)\subseteq \PNA(K)\cup\{0\}.
\]
Define \(T=EJ:c_0\to \Lip(M)\).  Since \(E\) is linear, \(T\) is linear.  For
every \(a\in c_0\),
\[
  \|T a\|_{\Lip(M)} \le \|E\|\,\|J a\|_{\Lip(K)}=\|a\|_\infty.
\]
On the other hand, restriction to \(K\) gives
\[
  \|T a\|_{\Lip(M)}
  \ge \|(T a)|_K\|_{\Lip(K)}
  = \|J a\|_{\Lip(K)}
  = \|a\|_\infty.
\]
Thus \(T\) is a linear isometry.

It remains to check pointwise norm-attainment.  Let \(a\in c_0\setminus\{0\}\).
Since \(J a\in \PNA(K)\), there exists \(p\in K\) such that
\[
  \sup_{q\in K\setminus\{p\}}
  \frac{|J a(p)-J a(q)|}{d(p,q)}
  =
  \|J a\|_{\Lip(K)}
  =
  \|a\|_\infty.
\]
Because \(T a\) extends \(J a\), the same slopes are present in \(M\):
\[
  \sup_{q\in M\setminus\{p\}}
  \frac{|T a(p)-T a(q)|}{d(p,q)}
  \ge
  \sup_{q\in K\setminus\{p\}}
  \frac{|J a(p)-J a(q)|}{d(p,q)}
  =
  \|a\|_\infty.
\]
The opposite inequality follows from \(\|T a\|_{\Lip(M)}=\|a\|_\infty\).
Hence \(T a\in\PNA(M)\).  Therefore \(T(c_0)\subseteq\PNA(M)\cup\{0\}\), as
claimed.
\end{proof}

\begin{corollary}[1-Lipschitz retract case]
\label{cor:retract}
Let \(K\subseteq M\) be pointed metric spaces.  Suppose there is a
1-Lipschitz retraction \(r:M\to K\) with \(r|_K=\operatorname{Id}_K\).  If
\(\PNA(K)\) contains an isometric copy of \(c_0\), then \(\PNA(M)\) contains
an isometric copy of \(c_0\).
\end{corollary}

\begin{proof}
Define \(E:\Lip(K)\to\Lip(M)\) by \(Ef=f\circ r\).  This operator is linear,
extends each \(f\), and satisfies \(\|E\|\le 1\) because \(r\) is
1-Lipschitz.  Since \(E\) is an extension operator, \(\|E\|=1\) unless
\(\Lip(K)=\{0\}\), and the trivial case is irrelevant for containing \(c_0\).
Theorem \ref{thm:lifting} applies.
\end{proof}

\section*{Consequences for the Source Problem}

Combining Corollary \ref{cor:retract} with the source's positive classes gives
ambient-space positive cases of Question 4.1.  For example, if \(M\) contains
an infinite pointed 1-Lipschitz retract \(K\) for which the source proves
\(c_0\subseteq\PNA(K)\) isometrically, then \(c_0\subseteq\PNA(M)\)
isometrically.  This includes retracts \(K\) that are not uniformly discrete
and the uniformly discrete classes covered by the source's Section 4
criteria.

\section*{Limitations}

This packet does not prove that every source subspace \(M_0\subseteq M\)
constructed in Theorem 4.10 admits such an extension operator.  The theorem
therefore leaves the full Question 4.1 open.  Its role is to isolate a sharp
and checkable extension hypothesis under which the subspace result upgrades
to the ambient space.

\section*{Novelty Check}

The local run indexes were searched for arXiv:2312.00393, pointwise
norm-attainment, \(c_0\), \(1\)-Lipschitz retracts, norm-one linear extension
operators, and extension-lifting language.  No duplicate packet or attempt was
found.  A bounded web search on June 29, 2026 for the same phrases found the
source paper but no later explicit answer or matching extension-lifting
statement.

\section*{Human Review Recommendation}

Review this as a partial permanence result, not as a full solution of Question
4.1.  The proof is short; the main review question is whether the norm-one
linear extension hypothesis is useful enough to count as a substantive
partial result toward the source's subspace-to-ambient obstruction.

\begin{thebibliography}{9}

\bibitem{CJLR231200393}
G. Choi, M. Jung, H. J. Lee, and O. Roldan,
\emph{Embeddings of infinite-dimensional spaces in the sets of norm-attaining
Lipschitz functions},
arXiv:2312.00393, 2023.

\end{thebibliography}

\end{document}
